% SOURCE_STATEMENT: "Can the quasi-bosonic approximation established for the first eigenvalue be extended to the true many-body ground state, possibly after suitable unitary transformations, and does it yield a quasi-bosonic description of the higher excitation spectrum?"
% Preparation and final self-audit: CONTRIBUTING.md prompts C and D; source/audit artifacts retained in the local physics_statements_2026_10_11 batch.
% !TeX program = lualatex
\documentclass[11pt]{article}
\usepackage[margin=25mm]{geometry}
\usepackage{fontspec}
\setmainfont{Latin Modern Roman}
\usepackage{amsmath,amssymb,mathtools}
\usepackage{unicode-math}
\setmathfont{Latin Modern Math}
\usepackage{xurl,hyperref}
\hypersetup{colorlinks=true,urlcolor=blue}
\setcounter{secnumdepth}{0}
\setlength{\parindent}{0pt}
\setlength{\parskip}{5pt}
\setlength{\emergencystretch}{3em}
\newcommand{\sref}[2]{\hyperlink{source:#1:#2}{[S#2]}}
\newcommand{\cataloguescope}[1]{\paragraph{Recorded target.}#1}
\newcommand{\R}{\mathbb R}
\newcommand{\C}{\mathbb C}
\newcommand{\N}{\mathbb N}
\newcommand{\Z}{\mathbb Z}
\newcommand{\Q}{\mathbb Q}
\newcommand{\F}{\mathbb F}
\newcommand{\eps}{\varepsilon}
\begin{document}
% UnsolvedMath ID 30006480; source catalogue code OWR-14299579-002
\setcounter{section}{30006479}
\section{Quasi-Bosonic Ground States and Excitation Spectra}\label{30006480}
{\small UnsolvedMath ID 30006480 · Mathematical Physics}
\subsection{Definitions and mathematical statement}

\paragraph{Shared notation.}$\mathbb N=\{1,2,\ldots\}$; the standard number fields have their usual meanings. All additional conventions are stated locally.


Bosonization approximates collective fermionic excitations by operators with approximately bosonic commutation relations. A quasi-bosonic energy approximation concerns an eigenvalue, whereas a ground-state approximation concerns a vector or density matrix and its observables. A unitary transformation can change the excitation representation without changing the spectrum.

\paragraph{Catalogue wording retained for traceability.}
Can the quasi-bosonic approximation established for the first eigenvalue be extended to the true many-body ground state, possibly after suitable unitary transformations, and does it yield a quasi-bosonic description of the higher excitation spectrum?

\paragraph{Scope and source qualification.}Energy convergence alone does not imply convergence of states. The statement retains both the transformed-state and excited-spectrum questions.
\subsection{Short English statement}
Extend the quasi-bosonic energy approximation to the actual many-body ground state and higher excitation spectrum.
\subsection{Sources}
\begin{itemize}
\item[\hypertarget{source:30006480:1}{S1}] ULAM.AI and the UnsolvedMath Contributors, UnsolvedMath, record 30006480 (OWR-14299579-002). Catalogue curation under CC BY 4.0; definitions and English exposition added locally.
\url{https://huggingface.co/datasets/ulamai/UnsolvedMath}.
\item[\hypertarget{source:30006480:2}{S2}] Arbeitsgemeinschaft: Analysis of Many-body Quantum Systems (2025), Oberwolfach Reports, DOI 10.4171/owr/2025/46, pp. 2491--2544. Original report and mathematical context. Incorporated report material remains under the source CC BY-SA 4.0 terms unless otherwise noted.
\url{https://doi.org/10.4171/owr/2025/46}.
\end{itemize}
\label{end:30006480}
\end{document}
